\documentclass[backaddress=true]{scrlttr2}
\usepackage[utf8x]{inputenc}
\usepackage{color,calc,mathptmx}
\usepackage{eurosym}
\usepackage{booktabs}
\usepackage[ngerman]{babel} 

\renewcommand{\familydefault}{\sfdefault}


\setkomavar{fromname}{Dipl.-Math. \\ Christian}
\setkomavar{fromaddress}{Kennich 22\\ Hamburg}
\setkomavar{subject}{Rechnung}
\setkomavar{place}{Hamburgo}
%\setkomavar{date}{31. Oktober 2014}


\firsthead{\null\hfill
  \parbox[t][\headheight][t]{5cm}{%
   \raggedright
   \color[gray]{.5}%
   \usekomavar{fromname}\\
   Sonninstraße 22\\
   20097 Hamburg\\[\baselineskip]
 %  \usekomavar*{fromphone}\\[\baselineskip]
   Steuernummer: 11/15111/11111\\[\baselineskip]
   }
}


\setkomavar{location}{
%\\[\baselineskip]
%\\[\baselineskip]
  \large{Rechnung  Nr. 1/2125
}}
\begin{document}
\date{18. Juli 2025}
\setlength{\parindent}{0pt}


\begin{letter}{
Mustermann Mann \\
Muster Str. 3c\\
12345 Musterstadt
}

\opening{}
%\renewcommand{\arraystretch}{1.3}
\begin{tabular}{|c|c|p{0.48\textwidth}|r|r|}
 \hline
  \textbf{Pos} & \textbf{Menge} & \textbf{Bezeichnung} & \textbf{Einzelpreis} & \textbf{Gesamtpreis} \\
 \hline
 1 & 1 Std &  Nachhilfe Statistik  \newline  &   \EUR{50,42} & \EUR{50,42} \\
 \hline

 %\\ \hline
 \multicolumn{4}{r}{\textbf{Zwischensumme}} & \multicolumn{1}{r}{\EUR{50,42}}\\
 \hline
 \multicolumn{4}{r}{\textbf{Gesamtbetrag (brutto) zuzüglich 19 Prozent Mehrwertsteuer}} & \multicolumn{1}{r}{\EUR{60}}\\

\end{tabular}
\\[\baselineskip]
%\textbf{Im ausgewiesenen Betrag ist gemäß § 19 UStG keine Umsatzsteuer enthalten.} %\\[\baselineskip]
%\textbf{Zahlbar bis 15.05.2012.}

%\vspace{35pt}
%\textbf{Zahlung abzüglich 3 Prozent Skonto bei Zahlung innerhalb von 14 Tagen }
%\\[\baselineskip]
%\\ 
\"Uberweisung auf IBAN DE12 1234 1234 1234 1234 12
%\\ Es wurde eine Anzahlung von \EUR{100} vereinbart.
%\closing{}
\vspace{35pt}
\\ \textbf{Dipl.-Math. Christian Heinrich Marek}
\end{letter}
\end{document}
